% !TeX root = ../../../manual.tex
\chapter{Institutions}

\section{Overview}

OmniLaTeX ships with 16 institution configurations located in
\texttt{config/institutions/}. Each institution package provides:
\begin{itemize}
    \item Logo definitions via \verb|\includesvg| or \verb|\includegraphics|
    \item Primary and secondary color definitions through \verb|\setkomacolor|
    \item Optional custom title pages built with TikZ overlays
    \item Translation-aware branding strings via \verb|\DeclareTranslation|
\end{itemize}

To activate an institution, load its package in your preamble or pass the
institution name through the OmniLaTeX configuration mechanism. The default
institution is \texttt{tuhh}, which serves as the reference implementation.

\section{Institution Reference Table}

\begin{longtblr}[
    caption = {Institution configurations shipped with OmniLaTeX},
    label = {tab:institutions},
]{colspec = {l S[table-format=3.0(3.0)] l l}, rowhead = 1}
    \toprule
    Institution & {Primary Color (RGB)} & {Has Logo} & {Custom Title} \\
    \midrule
    tuhh        & {---}                 & YES        & YES           \\
    cambridge   & (163,193,173)         & commented  & commented     \\
    cmu         & (196,18,48)           & commented  & commented     \\
    columbia    & (185,217,235)         & commented  & commented     \\
    epfl        & (255,0,0)             & commented  & commented     \\
    eth         & (31,64,122)           & commented  & commented     \\
    harvard     & (165,28,48)           & commented  & commented     \\
    imperial    & (0,62,116)            & commented  & commented     \\
    mit         & (163,31,52)           & commented  & commented     \\
    oxford      & (0,33,71)             & commented  & commented     \\
    princeton   & (231,117,0)           & commented  & commented     \\
    stanford    & (140,21,21)           & commented  & commented     \\
    tudelft     & (0,166,214)           & commented  & commented     \\
    tum         & (0,101,189)           & commented  & commented     \\
    yale        & (0,53,107)            & commented  & commented     \\
    generic     & {---}                 & commented  & commented     \\
    \bottomrule
\end{longtblr}

Institutions marked \emph{commented} ship with placeholder definitions that
can be activated by uncommenting the relevant sections and supplying the
appropriate logo files.

\section{TUHH in Detail}

The \texttt{tuhh} institution is the most complete configuration and serves as
the template for all others. Key features include:

\subsection{Logo Definitions}

The TUHH logo is defined as an SVG and embedded via \verb|\includesvg|:

\begin{verbatim}
\DeclareTranslation{english}{LogoInstitution}{%
    \includesvg[height=1.2cm]{tuhh-logo}%
}
\end{verbatim}

Bilingual support is achieved through separate translations for German and
English:

\begin{verbatim}
\DeclareTranslation{ngerman}{LogoInstitution}{%
    \includesvg[height=1.2cm]{tuhh-logo-de}%
}
\end{verbatim}

\subsection{Custom Title Style}

The \verb|\DefineTUHHTitleStyle| command creates a TikZ overlay title page
with the institution logo, a colored banner, and the document metadata:

\begin{verbatim}
\DefineTUHHTitleStyle{%
    \begin{tikzpicture}[remember picture, overlay]
        % Background fill
        \fill[tuhhblue] ...
        % Logo placement
        \node[anchor=north east] ...
        % Title text
        \node[anchor=west] ...
    \end{tikzpicture}%
}
\end{verbatim}

\subsection{TUHHuman TikZ Pic}

The TUHH configuration also defines a \verb|TUHHuman| TikZ \verb|pic| for
decorative use on title pages and section dividers. It is invoked with:

\begin{verbatim}
\pic at (0,0) {TUHHuman};
\end{verbatim}

\section{Creating Your Own Institution}

Follow these steps to add a new institution configuration:

\begin{enumerate}
    \item \textbf{Create the directory:}
\begin{verbatim}
mkdir config/institutions/myuni/
\end{verbatim}

    \item \textbf{Create the package file} \texttt{myuni.sty} with a proper
          package declaration:
\begin{verbatim}
\ProvidesPackage{myuni}[2026/01/01 My University Config]
\end{verbatim}

    \item \textbf{Define the logo} using \verb|\DeclareTranslation| so that
          it is language-aware:
\begin{verbatim}
\DeclareTranslation{english}{LogoMyuni}{%
    \includesvg[height=1.2cm]{myuni-logo}%
}
\end{verbatim}

    \item \textbf{Set institution colors} via KOMA-Script color commands:
\begin{verbatim}
\setkomacolor{title}{RGB}{0,50,100}
\setkomacolor{sectioning}{RGB}{0,50,100}
\end{verbatim}

    \item \textbf{Optionally define a custom title page} using TikZ overlays,
          following the TUHH pattern shown above.
\end{enumerate}

\section{Logo Management}

OmniLaTeX uses \verb|\includesvg| from the \texttt{svg} package to embed
vector graphics. The workflow requires:

\begin{itemize}
    \item \textbf{Inkscape} must be installed and available on \verb|PATH|.
          The \texttt{svg} package invokes Inkscape to convert SVG files to
          PDF on the fly.
    \item \textbf{\texttt{-{}-shell-escape}} must be passed to the LaTeX
          engine to allow external command execution.
    \item \textbf{\texttt{\textbackslash svgpath}} configures search paths for
          SVG files, analogous to \verb|\graphicspath|:
\begin{verbatim}
\svgpath{{images/logos/}{assets/}}
\end{verbatim}
    \item Converted PDFs are cached alongside the source SVG so subsequent
          compilations are faster.
\end{itemize}

For debugging SVG inclusion, consult the generated \texttt{.svg-inkscape}
log files in the output directory.
